\documentclass[11pt]{article}
\usepackage{amsmath, amssymb, amsthm, geometry}
\geometry{margin=1in}

\title{A Scalable Layer--Propagation and MRV Branching Method for Solving and Generating $N \times N$ Sudoku}
\author{Shaimaa Said Soltan}
\date{}

\begin{document}
\maketitle

\begin{abstract}
This paper presents a general Sudoku solving and puzzle-generation framework
capable of handling arbitrary grid sizes $N \times N$, including large orders such as
$12 \times 12$, $16 \times 16$, $25 \times 25$, and $36 \times 36$.
The method is based on a hybrid algorithm combining layer-based constraint propagation,
single-missing deduction, minimum-remaining-values (MRV) heuristics, shallow branching,
and early conflict detection. 
The solver is fully consistent with the implementation provided in the attached code,
and the generator uses symmetry-preserving and structure-aware clue removal to ensure
solvability even for large grids.
We analyze the novelty of this approach and compare its performance to classical
backtracking and human-style logical solvers.
\end{abstract}

\section{Introduction}
Sudoku solving algorithms traditionally fall into two categories:
\begin{itemize}
    \item \textbf{Logical solvers}, which apply human-style deduction rules such as
    naked singles, hidden singles, pairs, triples, and advanced pattern recognition.
    \item \textbf{Backtracking solvers}, which perform depth-first search over candidate
    assignments until a valid solution is found.
\end{itemize}

Logical solvers do not scale well to large $N$, because the number of patterns grows
combinatorially. Pure backtracking becomes infeasible for $N > 9$ due to the exponential
search space. The solver described in this paper, based directly on the attached code,
combines propagation and limited branching in a way that preserves both speed and
generality.

\section{Grid Structure and Box Shapes}
For any valid Sudoku size $N$, the grid is partitioned into subregions of size
$a \times b$, where $ab = N$. The implementation automatically selects the box shape
$(a,b)$ that is closest to a square by minimizing


\[
|a - \sqrt{N}|.
\]


Examples include:


\[
9 = 3 \cdot 3,\quad 12 = 3 \cdot 4,\quad 16 = 4 \cdot 4,\quad 36 = 6 \cdot 6.
\]


This generality allows the solver to support any $N$ for which a rectangular factorization exists.

\section{Layer-Based Propagation}
The core innovation of the solver is the \emph{layer construction} for each digit.
For a given digit $d$, a boolean layer matrix is built indicating which empty cells
\emph{could} contain $d$ after blocking:
\begin{itemize}
    \item all cells in any box already containing $d$,
    \item all cells in any row containing $d$,
    \item all cells in any column containing $d$.
\end{itemize}

If a box has exactly one remaining allowed position for $d$, the digit is placed.
This captures many classical logical rules in a single unified mechanism.

\section{Single-Missing Deduction}
After layer propagation, the solver performs deterministic filling of:
\begin{itemize}
    \item any row with exactly one missing value,
    \item any column with exactly one missing value.
\end{itemize}
This is a powerful deduction step that often resolves large portions of the grid.

\section{MRV Candidate Selection}
When propagation stalls, the solver computes candidate sets for every empty cell.
The \emph{minimum remaining values} (MRV) heuristic selects the cell with the fewest
candidates:


\[
\text{MRV}(r,c) = \arg\min_{(r,c)} |\text{Candidates}(r,c)|.
\]


This dramatically reduces branching and is essential for scaling to large $N$.

\section{Shallow Branching with Forbidden-Value Learning}
The solver performs only one level of branching:
\begin{enumerate}
    \item Select the MRV cell.
    \item Try each candidate value.
    \item After each attempt, run full layer propagation.
    \item If a contradiction occurs, mark the value as forbidden.
\end{enumerate}

If all candidates fail, the puzzle is declared unsolvable.
This shallow branching is far more efficient than full backtracking, yet powerful
enough to resolve puzzles that propagation alone cannot solve.

\section{Conflict Detection}
A contradiction is detected if:
\begin{itemize}
    \item a row contains duplicate digits,
    \item a column contains duplicate digits,
    \item a box contains duplicate digits,
    \item an empty cell has zero candidates.
\end{itemize}
These checks prune the search space aggressively and prevent wasted branching.

\section{Full-Grid Generation}
The generator constructs a valid Latin-based full solution using the formula:


\[
\text{grid}[r][c] = (r \cdot b + \lfloor r/a \rfloor + c) \bmod N + 1,
\]


where $(a,b)$ is the box shape. The grid is then randomized by:
\begin{itemize}
    \item permuting digits,
    \item shuffling rows within row bands,
    \item shuffling columns within column bands.
\end{itemize}

\section{Puzzle Generation for Large $N$}
Random clue removal fails for large grids such as $36 \times 36$.
The attached implementation uses:
\begin{itemize}
    \item \textbf{symmetry-preserving removal} (180° rotational pairs),
    \item \textbf{minimum clues per row, column, and box},
    \item \textbf{delayed solvability checks} (only near the target clue count),
    \item \textbf{the hybrid solver itself} to verify solvability.
\end{itemize}
This produces solvable puzzles efficiently even for $36 \times 36$.

\section{Performance Comparison}
We compare three solver types:

\subsection{Pure Backtracking}
Backtracking explores a search tree of size approximately


\[
O(N^{N^2}),
\]


which is infeasible for $N > 9$.

\subsection{Human-Style Logical Solvers}
These rely on pattern recognition and do not generalize well to large $N$.
Many rules become ineffective when candidate sets are large.

\subsection{Hybrid Layer--MRV Method (This Work)}
The attached solver reduces the search space by:
\begin{itemize}
    \item eliminating most candidates through layer propagation,
    \item selecting the most constrained cell via MRV,
    \item performing only shallow branching,
    \item pruning contradictions early.
\end{itemize}

Empirical performance from the implementation:
\begin{itemize}
    \item $9 \times 9$: milliseconds,
    \item $12 \times 12$: under a second,
    \item $16 \times 16$: seconds,
    \item $36 \times 36$: seconds to tens of seconds depending on clue density.
\end{itemize}

\section{Conclusion}
The solver and generator described here provide a practical and scalable approach
to Sudoku of arbitrary size. The combination of layer propagation, MRV heuristics,
shallow branching, and early conflict detection achieves a balance between speed
and generality that is not present in classical approaches. The generator's
structure-aware clue removal ensures solvability even for large grids such as
$36 \times 36$.

\end{document}
